\documentclass{report}
\usepackage[utf8]{inputenc}
\usepackage[T1]{fontenc}
\usepackage{geometry}
\usepackage{graphicx}
\usepackage{float}
\usepackage{hyperref}
\usepackage[english,indonesian]{babel}

\title{Optimasi Pemrosesan Template LaTeX Menggunakan Pendekatan Berbasis AI}
\author{Academic Support Agent}
\date{\today}

\begin{document}
\maketitle

\begin{abstract}

\end{abstract}

\chapter{PENDAHULUAN}


\chapter{TINJAUAN PUSTAKA}


% Original content (commented out):
% % Template do Relatorio das Atividades
% Disciplina GSI524 - Redes de Computadores
% Prof. Thiago Pirola Ribeiro
% Criado em: 2019/02/21

\documentclass[12pt,a4paper]{report}

\pagestyle{empty}

\usepackage[utf8]{inputenc}
\usepackage[brazil]{babel}
\usepackage{graphics}
\usepackage{graphicx,url}
\usepackage[top=2cm, bottom=2cm, left=2cm, right=2cm]{geometry}
\usepackage{enumitem} % para colocar (a)
\hyphenpenalty=5000 % para evitar que ultrapasse as margens e faca hifenacao
\tolerance=1000     % para evitar que ultrapasse as margens e faca hifenacao

\begin{document}
	\begin{center}
		
		{\Large Universidade Federal de Uberlândia  - UFU}
		
		Bacharelado em Sistemas de Informação - Campus Monte Carmelo
		
		GSI524 – Redes de Computadores – 2019/1
		
		\textbf{SEU NOME - SUA MATRÍCULA}
	
	\end{center}
	
	\vspace{5pt}

	\begin{center}
		{\LARGE \textbf{Relatório 1\\ 	\vspace{10pt} Meu Primeiro Relatório}}
	\end{center}
	
	\vspace{5pt}
	
	\begin{enumerate}
		\item Resposta da Primeira Questão:
		\begin{enumerate}[label=(\alph*)]
			\item É uma questão muito complexa.
			\item Algo fácil para muitos e mais fácil para mim.
		\end{enumerate}

		\begin{figure}[ht]
			\label{fig:logoUFU}
			\centering
			\includegraphics[scale=0.15]{imagens/ufu.jpeg}
			\caption{Logo da UFU.}
		\end{figure}

		\item A questão 2 tende a ser muito mais tranquila que a questão 723. Na Figura \ref{fig:logoUFU} observa-se o símbolo da UFU.
	\end{enumerate}
	
\end{document}

\end{document}